\subsection{Which data are common to the two continuations?}\label{sec:suppdom}
Start with one source preparation, an independent unknown qubit \(C\), initially empty registers, and the same scattering block and charge rule at each receiver. In the flat experiment the space of records is the four-bit set \((n_A,m_A,n_B,m_B)\). On this network an observer who receives only two bits sums over the other two. This partition changes conditional probabilities while preserving the complete channel, including in a network that chooses its next operation from earlier outcomes. At each reachable prefix the available operations obey \(\sum_z L_z^\dagger L_z=I\) on its input domain. Iterating this identity from the leaves gives normalization at the root.

To describe what happens to the qubit, specify in addition the observer's material port, the incoming signals that reach it by its clock deadline, and the observables available after the contact. In the spherical sector its inward fiber is \(D_O=H^0(\mathbb{CP}^1,{\cal O}(m_O-n_O))\otimes C\). Write \(S_O\) for the surface carrier. The local outgoing space is \(D_O\oplus S_O\). Consider the particular joint receipt \(n_A=m_A=n_B=1,\ m_B=0\).
\begin{center}\small
\begin{tabularx}{\linewidth}{@{}lXX@{}}\toprule
 & A at the matched port & B at the deficient port\\\midrule
Physical receipt & \(q_A=0\) before A's contact & \(q_B=-1\) before B's contact\\
Allowed inward fiber & One spatial mode carrying \(C\) & Zero spatial modes\\
Map on an arbitrary \(\psi\in C\) & \(V_A\psi=(\psi,0)\) in \(D_A\oplus S_A\) & \(V_B\psi=(0,J_B\psi)\), \(J_B:C\to S_B\) isometric\\
Accessible reading in this example & A may measure the inward code and its receipt & B may measure its receipt and the surface code or its specified radiation carriers\\
\bottomrule\end{tabularx}
\end{center}
The displayed maps are alternative continuations of the same input; multiplying them into a joint two-qubit output would change the experiment. On a different B-history the local inward mode can exist while remaining inaccessible to an exterior decoder; this produces the erasure term in the article. A zero physical fiber in the selected B-history is a stronger statement. It requires the assumed restriction on \emph{all} inward ports in the band.

There is a direct overlap test. Copy the already classical bit \(n_A\) into two ideal causal messages in a calibration of the flat receiver, and allow time for both messages to arrive. Each observer then reads the same bit with probability \(\Pr(n_A=1)=p=0.64\). Their own route readings can still differ. In general a shared effect \(E\) has representations \(j_A(E),j_B(E)\) whose pullbacks by the two continuation maps agree on the common input. This requirement concerns the part of the observable algebra that can actually be transported. It assumes no exchange between black-hole endpoints.

\paragraph{Changing a clock order.}
For two incomparable contacts the retained subchannels must agree in either allowed order, with physical labels, output spaces and subsequent controls transported along the same identification. In common coordinates the condition is
\(\mathcal I_{e,z}\circ\mathcal I_{f,w}=\mathcal I_{f,w}\circ\mathcal I_{e,z}\)
for each \((z,w)\). The separated local operators of the flat fixture act on different tensor factors and satisfy it even for an entangled source; causal predecessors keep their order. Commutation only after summing outcomes gives a weaker test. On a single qubit let each instrument measure \(Z\) or \(X\), retain the sign, and reset the state to \(I/2\). The inclusive maps coincide. On \(|0\rangle\) the probability of the two positive records is \(1/2\) in the \(Z,X\) order and \(1/4\) in the \(X,Z\) order. These are sequential contacts on one qubit. A general theory of dressed gravitational contacts must supply its own outcome-wise comparison.
